% Capítulo 6 - Conclusões (PT-PT; espelho de ../thesis/ch6/chapter6.tex)
% Tradução pura: labels/refs/citações/números/estrutura idênticos; só a língua muda.

\chapter{Conclusões}
\label{chap:conclusions}

Esta dissertação propôs-se desenhar, implementar e avaliar um sistema de alertas financeiros
explicável e reprodutível para investidores de retalho dos EUA, movido por dois gatilhos
independentes e entregue através do Telegram. Este capítulo resume o que foi alcançado face às
questões de investigação, revisita as contribuições, e enuncia as limitações e as direções que elas
abrem.

\section{Conclusões Principais}
\label{sec:concl_main}

O trabalho resume-se melhor voltando às quatro questões de investigação do
Capítulo~\ref{chap:introduction}, reunidas num relance na Figura~\ref{fig:rq_scorecard} e depois
tomadas uma a uma.

\begin{figure}[ht]
\centering
\small
\renewcommand{\arraystretch}{1.5}
\begin{tabular}{p{4.4cm} l p{5.0cm}}
\toprule
\tabhead{Questão de investigação} & \tabhead{Veredicto} & \tabhead{Um número} \\
\midrule
RQ1: detetar movimentos abruptos de forma transparente
  & \colorbox{green!20}{\textbf{Sim}}
  & taxa de disparo $0.015$ vs $0.344$; $F_1$ $0.516$ \\
RQ2: precedentes análogos, sem lookahead
  & \colorbox{green!20}{\textbf{Sim} (recuperação)}
  & P@5 $0.514$ vs $0.346$ lexical, $0.240$ aleatório \\
RQ3: explicações fiéis e úteis
  & \colorbox{yellow!25}{Fiel; utilidade em aberto}
  & fiel por construção; ainda sem estudo humano \\
RQ4: triagem para além da volatilidade
  & \colorbox{yellow!25}{Mecanismo sim, texto não}
  & quota $0.163 \to 0.632$; PR-AUC $0.542$ vs $0.496$ \\
\bottomrule
\end{tabular}
\caption[As quatro questões de investigação num relance]{As quatro questões de investigação e os seus
veredictos, reportados exatamente como se revelaram. Duas são um sim sem reservas; duas são divisões
honestas: as explicações são fiéis por construção mas a sua utilidade aguarda um estudo humano, e a
triagem aprendida ajuda como mecanismo de produto embora nenhum modelo de texto de manchetes tenha
batido a linha de base da volatilidade neste corpus. Os números são do
Capítulo~\ref{chap:casestudies}.}
\label{fig:rq_scorecard}
\end{figure}

\paragraph{RQ1: deteção transparente de movimentos abruptos.}
\textbf{Sim.} Um detetor de \emph{z}-score deslizante assinala retornos anormais e, mais importante,
fá-lo a um ritmo estável \emph{em ações de volatilidade muito diferente}, onde um limiar de
percentagem fixa não consegue. Como reportado no
Capítulo~\ref{chap:casestudies}, a taxa de disparo variou apenas $0.015$ entre os quinze tickers no
\emph{z}-score, contra $0.344$ na linha de base fixa, e o \emph{z}-score atingiu um $F_1$ de $0.516$
contra um indicador de movimentos extremos, mais do dobro da linha de base. O método é simples o
suficiente para que cada alerta possa ser explicado ao investidor.

\paragraph{RQ2: precedentes análogos sem lookahead.}
\textbf{Sim, para a recuperação.} A recuperação semântica com Sentence-BERT recuperou muito mais
precedentes análogos por setor do que qualquer alternativa trivial: uma precisão@$5$ cross-ticker de
$0.514 \pm 0.015$ (em cinco corridas), contra $0.346$ de uma linha de base lexical, $0.240$ do
aleatório e $0.126$ da recência. Uma decomposição por setor mostrou que o ganho é maior onde o
vocabulário do domínio é mais distintivo (energia, saúde). O estudo de caso ponta a ponta mostrou-a a
trazer à superfície precedentes genuinamente temáticos a partir de fontes de empresas diferentes. O
impacto é medido estritamente após o evento, ao estilo de um estudo de evento, pelo que nenhuma
informação futura se infiltra na evidência. Isto foi depois validado à escala: no corpus FNSPID
multi-ano (cerca de oitenta mil manchetes) o mesmo protocolo cross-ticker atingiu uma precision@$5$
de $0.595$, acima do valor preliminar, e a propriedade tema-não-direção foi quantificada diretamente,
já que a consistência de direção dos clusters recuperados ($0.71$) fica logo acima do seu chão do
acaso de $0.69$.

\paragraph{RQ3: explicações fiéis e úteis.}
\textbf{Fiéis, sim; úteis, ainda em aberto.} Como cada alerta é composto diretamente a partir dos
mesmos objetos que o sistema calcula (o \emph{z}-score e os seus parâmetros, ou os precedentes
recuperados com as suas pontuações e impactos medidos), a explicação não pode divergir da lógica
subjacente, e cada alerta termina com um aviso explícito de não-previsão. A utilidade para um
investidor de retalho é plausível por desenho mas ainda não foi medida com um estudo humano, sendo
por isso reportada como um ponto em aberto e não como um resultado assente.

\paragraph{RQ4: triagem aprendida para além das linhas de base de volatilidade.}
\textbf{Não na hipótese do texto; sim no mecanismo.} Um modelo de materialidade supervisionado foi
treinado, calibrado e testado em $79{,}753$ exemplos do FNSPID (2018--2023) sob um protocolo temporal
estrito. Como mecanismo de produto, a triagem aprendida é claramente valiosa: dentro de um orçamento
de cinco alertas por dia, elevou a quota de alertas materiais de $0.163$ para $0.632$, com
probabilidades calibradas. Mas a comparação decisiva pré-comprometida foi noutro sentido: nenhum
modelo que lê o texto da manchete bateu a linha de base só-volatilidade (PR-AUC $0.542$ contra $0.496$
do modelo principal contexto$+$texto), pelo que neste corpus o sinal da triagem vive no contexto de
mercado, não na redação da manchete. Juntamente com a comparação com a Isolation Forest do estudo de
anomalias, este é o segundo teste justo e causal em que a escolha transparente venceu, e ambos os
resultados são reportados exatamente como se revelaram. O resultado negativo no texto foi ainda
sujeito a stress: um bootstrap por cluster (ticker, dia) confirma que adicionar o texto da manchete
baixa a precisão de forma robusta e não por acaso, e um re-teste deliberadamente justo, com
regularização afinada, um bloco de texto reduzido em dimensão e um encoder de domínio financeiro,
recupera o texto apenas até ao nível do contexto de mercado, nunca acima. O resultado é robusto, não
um artefacto de um pipeline de texto sub-ajustado.

\section{Questões de Investigação e Objetivos Alcançados}
\label{sec:concl_objectives}

O objetivo geral, desenhar, implementar e avaliar um sistema de alertas explicável e reprodutível
movido por dois gatilhos, foi cumprido: o InvestiGator é um sistema funcional cujos dois gatilhos
foram provados de ponta a ponta e cujos componentes centrais, incluindo o modelo de triagem treinado,
foram avaliados sobre dados reais. Os objetivos de apoio também foram cumpridos:
uma fatia fina de ponta a ponta foi entregue primeiro; cada componente foi mantido na sua forma mais
simples e defensável; a avaliação seguiu um desenho fixado à partida; e todo o sistema é reprodutível
a partir de scripts versionados.

\section{Contribuições Revisitadas}
\label{sec:concl_contributions}

Enquanto dissertação de \emph{Engenharia} de Inteligência Artificial, a contribuição é a integração,
aplicação e avaliação crítica de componentes existentes num todo funcional, explicável e reprodutível,
e não um novo algoritmo. Destacam-se quatro contribuições concretas. Primeira, uma metodologia
documentada e reprodutível para \textbf{correlação notícia--mercado} que combina recuperação por
embeddings de frase com janelas de impacto de estudo de evento, com escolhas explícitas de métrica de
similaridade e horizonte e evitação explícita de lookahead. Segunda, o \textbf{InvestiGator}, um
\textbf{sistema de alertas XAI-first} cuja cadeia de raciocínio inteira (deteção, evidência, fontes e
precedentes) é exposta a um utilizador não-especialista e é fiel por construção. Terceira, um
\textbf{modelo de triagem de materialidade} treinado e calibrado, rotulado pelo próprio código de
estudo de evento do sistema sobre o corpus multi-ano, integrado nos alertas desligado por defeito como
evidência ordenada e explicada, e continuamente pós-validado em produção contra os resultados que
efetivamente se seguiram. Quarta, uma \textbf{avaliação crítica e honesta} de cada componente central
contra linhas de base triviais, lexicais e de volatilidade sobre dados reais, incluindo duas
comparações pré-comprometidas estatística-versus-aprendida que os métodos transparentes venceram; os
resultados, ablações e limitações são todos reportados com clareza e regenerados a partir de scripts
versionados.

\section{Limitações}
\label{sec:concl_limitations}

Restam várias limitações, enunciadas com clareza:
\begin{itemize}
  \item \textbf{Corpus.} A avaliação de recuperação primária usou um corpus de notícias recente, de
        poucos meses, de um serviço gratuito, o que torna esses números preliminares; um seguimento no
        corpus FNSPID multi-ano validou a recuperação à escala (precision@$5$ de $0.595$). O estudo de triagem usou de facto o subconjunto
        multi-ano do FNSPID, mas com catorze dos quinze tickers (a Meta está indexada sob o seu
        símbolo antigo no corpus) e apenas o texto da manchete.
  \item \textbf{Indicadores aproximados.} A relevância é aproximada por concordância de setor, e o
        rótulo de anomalia é relativo à volatilidade, razão pela qual o argumento da taxa de disparo
        sem rótulos tem primazia.
  \item \textbf{Texto curto.} As notícias são representadas apenas por manchetes, o que limita a
        semântica captada.
  \item \textbf{Tema, não direção.} A similaridade por embeddings capta o tema, não a direção, pelo
        que um agrupamento recuperado pode misturar movimentos de subida e de descida; o impacto médio
        é evidência ao nível do tema, não um sinal direcional, razão pela qual os precedentes são
        sempre mostrados individualmente.
  \item \textbf{Sem estudo humano.} A utilidade para um investidor real ainda não foi medida.
  \item \textbf{Por desenho.} O sistema não faz previsão de preços nem negociação, pelo que as
        questões baseadas em lucro estão fora de âmbito, não por responder.
\end{itemize}

\section{Trabalho Futuro}
\label{sec:concl_future}

As limitações mapeiam diretamente em trabalho futuro, resumido na
Figura~\ref{fig:limitations_future}.

\begin{figure}[ht]
\centering
\begin{tikzpicture}[
  font=\footnotesize, >={Stealth[length=2.6mm]},
  lim/.style={draw, rounded corners, align=left, text width=4.5cm, inner sep=4pt,
              fill=orange!8, draw=orange!55, minimum height=1.05cm, anchor=west},
  fut/.style={draw, rounded corners, align=left, text width=6.2cm, inner sep=4pt,
              fill=green!9, draw=green!50!black, minimum height=1.05cm, anchor=west},
]
  \node[font=\small\bfseries, text=orange!65!black, anchor=west] at (0.1,0.95) {Limitação};
  \node[font=\small\bfseries, text=green!45!black, anchor=west] at (5.4,0.95) {Passo concreto seguinte};
  \node[lim] (l0) at (0,0) {Corpus recente, de poucos meses};
  \node[fut] (f0) at (5.3,0) {Estudo das magnitudes de impacto à escala (recuperação já validada no FNSPID)};
  \node[lim] (l1) at (0,-1.45) {Só manchetes (texto curto)};
  \node[fut] (f1) at (5.3,-1.45) {Texto mais rico: corpo dos artigos e uma feature de sentimento};
  \node[lim] (l2) at (0,-2.9) {Sem estudo humano de utilidade};
  \node[fut] (f2) at (5.3,-2.9) {Pequeno estudo com rubrica: clareza, completude, acionabilidade (RQ3)};
  \node[lim] (l3) at (0,-4.35) {Tema, não direção};
  \node[fut] (f3) at (5.3,-4.35) {De-duplicar precedentes e sinalizar divergência de direção};
  \draw[->, thick, black!55] (l0.east) -- (f0.west);
  \draw[->, thick, black!55] (l1.east) -- (f1.west);
  \draw[->, thick, black!55] (l2.east) -- (f2.west);
  \draw[->, thick, black!55] (l3.east) -- (f3.west);
\end{tikzpicture}
\caption[As limitações mapeiam em trabalho futuro]{Cada limitação enunciada mapeia diretamente num
passo concreto seguinte. Várias já viram uma primeira iteração do lado da produção (uma base de
conhecimento viva maturada contra preços observados, sinalização explícita de divergência de direção,
e uma cotação intradiária durante o horário de mercado), reportadas como seguimento de engenharia cuja
avaliação formal fica para trabalho futuro.}
\label{fig:limitations_future}
\end{figure}

\begin{itemize}
  \item Avaliar as \emph{magnitudes de impacto} no corpus FNSPID multi-ano. A própria recuperação já
        foi validada aí (precision@$5$ de $0.595$, com estatísticas de dispersão e de consistência de
        direção); o que resta é o estudo das magnitudes de impacto ajustadas ao mercado sobre esses
        precedentes multi-ano.
  \item Realizar um pequeno estudo humano, aplicando uma rubrica de clareza, completude e
        acionabilidade a um conjunto de alertas, para resolver a questão em aberto da utilidade (RQ3).
  \item Atacar a lacuna de texto em aberto da RQ4: texto mais rico do que manchetes (corpo dos
        artigos, sentimento financeiro como feature), o ticker em falta recuperado sob o seu símbolo
        antigo, e o ciclo de pós-validação ao vivo a acumular decisões rotuladas frescas para
        re-treino periódico; um tratamento por contextual-bandit do limiar do alerta é a extensão
        aprendida natural.
  \item Como produto, de-duplicar precedentes quase idênticos e sinalizar a divergência de direção
        dentro de um agrupamento recuperado, complementando a ordenação por severidade que a triagem já
        fornece contra os riscos de fadiga de alertas e de confiança excessiva do
        Capítulo~\ref{chap:investigator}.
\end{itemize}
A produção ao vivo, depois de a avaliação ter sido congelada, já motivou uma primeira iteração de
vários destes pontos: o sistema em produção faz agora crescer uma base de conhecimento \emph{viva} a
partir das manchetes que varre (cada caso maturado contra preços observados dias depois), ordena os
precedentes com um fator de decaimento por idade mostrando sempre a similaridade verdadeira e a idade
de cada caso, sinaliza explicitamente a divergência de direção, e avalia o movimento do dia em curso a
partir de uma cotação em tempo real durante o horário de mercado. A produção também ensinou a sua
própria lição, reportada no Capítulo~\ref{chap:investigator}: uma única fonte de preços gratuita
provou ser um ponto único de falha em infraestrutura alojada, e foi substituída por uma cadeia de
fallback multi-fornecedor. Estas iterações do lado da produção reutilizam os componentes validados sem
alterações e são aqui reportadas como seguimento de engenharia; a sua avaliação formal fica para
trabalho futuro. Um ponto está já \emph{validado} em vez de especulativo: a comparação estendida do
Estudo de Caso~1 mostrou que uma estimativa de volatilidade com ponderação exponencial melhora
substancialmente a concordância do detetor com o indicador, pelo que adotá-la é uma mudança medida, de
uma linha, à espera apenas de uma decisão de explicabilidade. Em conjunto, estas mudanças transformariam
um protótipo sólido num sistema mais exaustivamente validado, sem abdicar da simplicidade e
transparência que o tornam defensável.
